\documentclass[12pt]{article}

\setlength{\parindent}{0pt}
\setcounter{secnumdepth}{3}
\usepackage[margin=1in]{geometry}

% packages
\usepackage{amsmath, amsfonts, amsthm, amssymb}
\usepackage{graphicx, float}
\usepackage{mathtools}
\usepackage{titlesec}
\usepackage{interval}
\usepackage{hyperref}
\usepackage{siunitx}
\usepackage{titling}
\usepackage{vwcol}
\usepackage{setspace}
\usepackage{empheq}
\usepackage{cancel}
\usepackage{esdiff}
\usepackage{multicol}
\usepackage{mdframed}
\usepackage{esdiff}
\usepackage{tikzsymbols}
\usepackage{multicol}
\usepackage{tikz}
\usepackage{varwidth}
\usepackage{pgfplots}
\usepackage{fancyhdr}
\usepackage{enumitem}
\usepackage{tcolorbox}

\intervalconfig {
	soft open fences
}

% header
\pagestyle{fancy}
\fancyhf{}
\lhead{Ethan Fan}
\rhead{AP Calculus BC}
\cfoot{\thepage}

% section format
\titleformat{\section}
  {\large \bfseries}
  {}
  {0em}
  {}
\titlespacing*{\section}{0pt}{0.5em}{0.3em}

\titleformat{\subsection}[runin]
  {\bfseries}
  {}
  {0em}
  {}
\titlespacing*{\subsection}{0pt}{0pt}{0em}

\titleformat{\subsubsection}[runin]
  {\itshape}
  {(\roman{subsubsection})}
  {0em}
  {}

% commands
\newcommand{\alignedintertext}[1]{%
  \noalign{%
    \vskip\belowdisplayshortskip
    \vtop{\hsize=\linewidth#1\par
    \expandafter}%
    \expandafter\prevdepth\the\prevdepth
  }%
}

\newcommand*{\paren}[1]{\ensuremath\left(#1\right)}
\newcommand*{\limit}[2][x]{\ensuremath{\displaystyle\lim_{#1\to#2}}}
\newcommand*{\Limit}[3][x]{\ensuremath{\displaystyle\lim_{#1\to#2}\left[#3\right]}}
\newcommand*{\deriv}[1][x]{\ensuremath{\dfrac{\mathrm{d}}{\mathrm{d}#1}}}
\newcommand*{\Deriv}[2][x]{\ensuremath{\dfrac{\mathrm{d}}{\mathrm{d}#1}\left[#2\right]}}
\newcommand{\ddx}{\dfrac{d}{dx}}
\newcommand{\dydx}{\dfrac{dy}{dx}}
\newcommand{\dydt}{\dfrac{dy}{dt}}
\newcommand{\dxdt}{\dfrac{dx}{dt}}
\newcommand*{\iinteg}[2][x]{\ensuremath{\displaystyle\int #2\;\mathrm{d}#1}}
\newcommand*{\dinteg}[4][x]{\ensuremath{\displaystyle\int_{#2}^{#3}#4\;\mathrm{d}#1}}
\newcommand*{\abs}[1]{\ensuremath{\left|#1\right|}}
\newcommand*{\eps}{\varepsilon}
\newcommand*{\real}{\ensuremath\mathbb{R}}
\newcommand*{\integer}{\ensuremath\mathbb{Z}}
\newcommand*{\posint}{\ensuremath\mathbb{Z^+}}
\newcommand*{\cbrt}[1]{\ensuremath{\sqrt[3]{#1}}}
\newcommand*{\lheqzero}{\ensuremath{\underset{\text{L'H}}{\overset{\left[\frac00\right]}{=}}}}
\newcommand*{\lheqinfty}{\ensuremath{\underset{\text{L'H}}{\overset{\left[\frac{\infty}{\infty}\right]}{=}}}}
\newcommand{\tor}{\text{ or }}
\newcommand{\floor}[1]{\lfloor #1 \rfloor}

\newcommand{\vem}{\vspace{0.5em}}
\newcommand{\example}[1]{\section{Example #1}}
\newcommand{\problem}[1]{\section{Problem #1}}
\newcommand{\subproblem}[1]{\subsection{(#1)} \,}

\DeclareMathOperator{\dne}{DNE}
\DeclareMathOperator{\sgn}{sgn}

\newcommand{\definition}[1]{\begin{tcolorbox}[colback=red!5!white,colframe=red!75!black,parbox=false] #1 \end{tcolorbox}}

\newcommand{\theorem}[2]{\begin{tcolorbox}[title={#1},colback=blue!5!white,colframe=blue!75!black,parbox=false] #2 \end{tcolorbox}}

\allowdisplaybreaks
% \postdisplaypenalty=100000

% document
\begin{document}

\vspace*{0cm}

% opening
\begin{center}
    {\Large \bfseries Problem Set 15} \\[0.5cm]
    {\large Products of Powers of Trigonometric Functions} \\[0.35cm]
    {\normalsize October 15, 2026}
\end{center}

% problems

\problem{1}

\subproblem{a}
\begin{gather*}
    \int \sin^3 x\cos^2 x\, dx = \int (\cos^2 x-\cos ^4 x)\sin x\, dx\\
    =\int u^4-u^2\, du= \frac{u^5}{5}-\frac{u^3}{3}+C=\boxed{\frac{1}{5}\cos^5 x -\frac{1}{3}\cos^3x+C}
\end{gather*}

\subproblem{b}
\begin{gather*}
    u=3x \implies \frac{1}{3}du=dx\\
    \int \sin^43x\cos^3 3x\, dx = \frac{1}{3} \int \sin^4 u \cos^3 u\, du\\
    = \frac{1}{3} \int (\sin^4 u- \sin^6u)\cos u \, du\\
    =\frac{1}{3} \int v^4-v^6 \, dv = \frac{1}{3}(\frac{v^5}{5} - \frac{v^7}{7})+C\\
    = \boxed{\frac{1}{15} \sin^5 3x - \frac{1}{21} \sin^7 3x+C}
\end{gather*}

\subproblem{c}
\begin{gather*}
    \int \sin^2 5x\cos^2 5x\, dx\\
    =\frac{1}{4}\int (1-\cos 10x)(1+\cos 10x)\, dx\\
    =\frac{1}{4}\int 1-\cos^210x\, dx\\
    =\frac{1}{4} \int 1-\frac{1}{2} - \frac{\cos 20x}{2}\, dx\\
    =\frac{1}{8}(x-\frac{\sin 20x}{20})+C\\
    =\boxed{\frac{1}{8}x - \frac{1}{160} \sin 20x+C}
\end{gather*}

\subproblem{d}
\begin{gather*}
    \int \cos^5x \, dx = \int (1-2\sin^2x +\sin^4x)\cos x\, dx\\
    =\int u^4-2u^2+1\, du = \boxed{\frac{1}{5}\sin^5x -\frac{2}{3}\sin^3 x+\sin x+C}
\end{gather*}

\problem{2}

\subproblem{a}
\begin{gather*}
    f_{ave} = \frac{\int_{-\pi}^\pi \sin^2 x \cos^3 x\, dx}{2\pi }\\
    \int_{-\pi}^\pi \sin^2 x \cos^3 x\, dx =  2 \int_{0}^\pi \sin^2 x \cos^3 x\, dx\\
    = 2 \int_{0}^\pi (\sin^2 x - \sin^4 x) \cos x\, dx = 2 \int_{0}^{0} u^2 - u^4  \, du = 0\\
    f_{ave} = \boxed{0}\\
    u =  \pi -x \implies \int_0^{\pi} \sin^2 u \cos^3 u\, du =  -\int_{-\pi}^0 \sin^2 u \cos^3 u\, du
\end{gather*}

\subproblem{b}
\begin{gather*}
    \int_{-\pi}^\pi \sin^2 x \cos^2  x \, dx = 2 \int_0^\pi \sin^2x\cos^2x\,dx\\
    = \frac{1}{2} \int_0^\pi (1-\cos 2x)(1+\cos 2x)\,dx\\
    = \frac{1}{2} \int_0^\pi 1-\cos^2 2x\,dx\\
    = \frac{1}{2} \int_0^\pi 1-\frac{1}{2}-\frac{1}{2} \cos 4x\,dx\\
    = \frac{1}{2}(\frac{1}{2}x - \frac{1}{8}\sin 4x) \Big|^\pi_0 = \frac{\pi}{4}\\
    g_{ave} = \frac{\pi}{4}\cdot \frac{1}{2\pi} = \boxed{\frac{1}{8}}
\end{gather*}

\problem{3}

\subproblem{a}
\begin{align*}
    s(t) &= \int \sin (\omega t)\cos^2(\omega t)\, dt\\
    &= \begin{bmatrix}
        u =\cos \omega t\\
        du= -\omega \sin \omega t\, dt
    \end{bmatrix} 
    -\frac{1}{\omega} \int u^2\, du\\
    &= \boxed{- \frac{\cos^3(\omega t)}{3\omega } + \frac{1}{3\omega}}
\end{align*}

\subproblem{b}
\begin{gather*}
    - \frac{\cos^3(\omega t)}{3\omega } \in [-\frac{1}{3\omega},\frac{1}{3\omega}]\\
    \iff - \frac{\cos^3(\omega t)}{3\omega } + \frac{1}{3\omega} \in [0,\frac{2}{3\omega}]\\
    \cos^3(\omega t)= -1 \implies \boxed{t=\frac{\pi}{\omega}}
\end{gather*}

\subproblem{c}
Yes, it is periodic.
\begin{gather*}
    t= \frac{2\pi}{\omega } \\
    s(\frac{\pi}{\omega}) - s(0)= \frac{2}{3\omega}\\
    2\cdot \frac{2}{3\omega} = \boxed{\frac{4}{3\omega}}
\end{gather*}

\problem{4}

\subproblem{a}
\begin{gather*}
    \int \tan^3 x \sec^3 x \, dx = \int \sec^2x(\sec^2x -1) \sec x \tan x \, dx\\
    =\int u^4-u^2\, du = \frac{u^5}{5} -\frac{u^3}{3}+C\\
    =\boxed{\frac{1}{5} \sec^5 x-\frac{1}{3}\sec^3 x+C}
\end{gather*}

\subproblem{b}
\begin{gather*}
    \int \tan^4 x \sec^4x \, dx = \int \tan^4x(1+\tan^2x) \sec^2 x\, dx\\
    =\int u^4+u^6\,  du = \boxed{\frac{1}{7}\tan^7x+\frac{1}{5}\tan^5x+C}
\end{gather*}

\subproblem{c}
\begin{gather*}
    \int \tan^2 x\sec x\, dx = \int \sec^3 x-\sec x \, dx\\
    =\boxed{ \frac{1}{2}\sec x\tan x -\frac{1}{2}\ln |\sec x + \tan x |+C}
\end{gather*}

\problem{5}

\subproblem{a}
\begin{align*}
    \int \sec^nx\, dx &= \begin{bmatrix}
        u=\sec^{n-2}x & dv=\sec^2 x\\
        du=(n-2)\sec^{n-2} x \tan x \, dx & v = \tan x 
    \end{bmatrix}\\
    & \,\,\,\,\,\,\,\sec^{n-2}x\tan x - (n-2)\int \sec^{n-2}x\tan^2 x\, dx\\
    &=  \sec^{n-2}x\tan x - (n-2)\int \sec^{n}x\, dx +(n-2)\int \sec^{n-2}x\, dx\\
    &= \boxed{\frac{\sec^{n-2}x\tan x}{n-1} +\frac{n-2}{n-1}\int \sec^{n-2}x\, dx}
\end{align*}

\subproblem{b}
\begin{gather*}
    \int \sec^3 x \, dx = \boxed{\frac{1}{2}\sec x \tan x + \frac{1}{2} \ln |\sec x +\tan x |+C}
\end{gather*}

\subproblem{c}
\begin{gather*}
    \int \sec^5  x\,d x = \frac{\sec^3x\tan x}{4} + \frac{3}{4}(\frac{1}{2}\sec x \tan x + \frac{1}{2} \ln |\sec x +\tan x |)+C\\
    =\boxed{\frac{1}{4} \sec^3 x\tan x +\frac{3}{8} \sec x \tan x +\frac{3}{8} \ln|\sec x +\tan x| +C}
\end{gather*}

\subproblem{d}
\begin{gather*}
    \int \tan^2 x \sec^3 x \,d x = \int \sec^5x-\sec^3 x \, dx\\
    =\boxed{\frac{1}{4} \sec^3 x\tan x -\frac{1}{8} \sec x \tan x -\frac{1}{8} \ln|\sec x +\tan x| +C}
\end{gather*}

\subproblem{e}
\begin{gather*}
    \int \sec^4 x\, dx= \boxed{\frac{1}{3}\sec^2 x\tan x +\frac{2}{3} \tan x +C}
\end{gather*}

\problem{6}

\subproblem{a}
\begin{gather*}
    \int \frac{\cos^3x}{\sqrt{\sin x }} \, dx =\int \frac{1-\sin^2x}{\sqrt{\sin x}} \, \cos x \, dx\\
    = \int \frac{1}{\sqrt{u}} - u^\frac{3}{2} \, du = 2\sqrt u - \frac{2}{5} u^\frac{5}{2}+C\\
    = \boxed{2\sqrt{\sin x} - \frac{2}{5}\sin ^\frac{5}{2}x+C}
\end{gather*}

\subproblem{b}
\begin{gather*}
    \int \frac{\sin^2x}{\cos^4 x} \, dx = \int \sec^2 x\tan ^2 x\, dx\\
    =\boxed{\frac{1}{3}\tan^3  x + C}
\end{gather*}

\subproblem{c}
\begin{gather*}
    \int \cot^3x \csc ^3 x \, dx = \int (\csc ^4 x - \csc ^2 x )\cot x \csc x \, dx\\
    = \boxed{\frac{1}{3}\csc^3 x -\frac{1}{5}\csc ^5 x +C}
\end{gather*}

\subproblem{d}
\begin{gather*}
    \int \tan^3 x\, dx = \int \frac{\sec^2x-1}{\sec x} \sec x\tan x \, dx\\
    = \boxed{\frac{1}{2}\sec^2 x - \ln |\sec x|+C}
\end{gather*}

\problem{7}

\subproblem{a}
\begin{gather*}
    \int_{-\pi} ^\pi \cos^2 (nx)\, dx = 2 \int _0^\pi \cos^2 nx\, dx\\
    = \int_0^\pi 1+ \cos 2nx\, dx =x+ \frac{1}{2n} \sin 2nx\Big|^\pi_0 = \boxed{\pi}
\end{gather*}

\subproblem{b}
\begin{gather*}
    E(a)=\int_{-\pi}^\pi \paren{f(x)-a\cos(nx)}^2\, dx\\
    = \int_{-\pi}^\pi f(x)^2-2a\cos (nx)f(x)+a^2\cos(nx)^2 \, dx\\
    = a^2 \int_{-\pi}^\pi \cos(nx)^2 \, dx - 2a\int_{-\pi}^\pi \cos (nx)f(x)\, dx + \int_{-\pi}^\pi f(x)^2\, dx\\
    = \pi a^2  - 2a\int_{-\pi}^\pi \cos (nx)f(x)\, dx + \int_{-\pi}^\pi f(x)^2\, dx\\
    a= \boxed{\frac{1}{\pi} \int_{-\pi}^\pi \cos (nx)f(x)\, dx }
\end{gather*}

\subproblem{c}
\begin{gather*}
    a_1=\boxed{-4}\\
    a_2=\boxed{1}
\end{gather*}

\subproblem{d}
\begin{gather*}
    a=\frac{1}{\pi} \int_{-\pi}^\pi x^2\cos (nx)\, dx= \frac{2}{\pi} \int_{0}^\pi x^2\cos (nx)\, dx\\
    = \begin{bmatrix}
        u=x^2 & dv=\cos nx \, dx\\
        du =2x\, dx & v =\frac{1}{n}\sin nx \, dx
    \end{bmatrix}
    \frac{2}{n\pi} x^2\sin nx \Big|^\pi_0 - \frac{4}{n\pi} \int_0^\pi x \sin nx \, dx\\
    = \begin{bmatrix}
        u=x & dv=\sin nx\, dx\\
        du=dx & v= -\frac{1}{n} \cos nx
    \end{bmatrix}
    +\frac{4}{n^2\pi} x\cos nx \big|^\pi_0 - \frac{4}{n^2\pi} \int_0^\pi \cos nx\, dx\\
    = \boxed{\frac{4(-1)^n}{n^2} }
\end{gather*}
The equation is a great approximation of the function $y=x^2$ with cosine.

\problem{8}

\subproblem{a}
\begin{gather*}
    I_n= \int_0^{\frac{\pi}{2}} \cos^n x\cos (nx)\, dx\\
    =\begin{bmatrix}
        u=\cos^n x & dv=\cos (nx)\, dx\\
        du=-n\cos^{n-1}x\sin x  & u=\frac{1}{n}\sin (nx)
    \end{bmatrix} \frac{1}{n}\cos^n x\sin nx\Big|^\frac{\pi}{2}_0 \\
    + \int_0^\frac{\pi}{2} \cos ^{n-1}x\sin x \sin (nx)\, dx = \boxed{\int_0^\frac{\pi}{2} \cos ^{n-1}x\sin x \sin (nx)\, dx}
\end{gather*}

\subproblem{b}
\begin{gather*}
    \cos ((n-1)x)=\cos nx \cos x +\sin nx\sin x\\
    2I_n=\int_0^{\frac{\pi}{2}} \cos^n x\cos (nx)\, dx + \int_0^\frac{\pi}{2} \cos ^{n-1}x\sin x \sin (nx)\, dx\\
    = \int_0^\frac{\pi}{2} \cos^{n-1}x(\cos nx \cos x +\sin nx\sin x)\, dx\\
    =\int_0^\frac{\pi}{2} \cos^{n-1}x\cos ((n-1)x)\, dx = I_{n-1}\\
    I_n = \frac{1}{2}I_{n-1}, I_0 = \frac{\pi}{2} \implies \boxed{I_n = \frac{\pi}{2^{n+1}}}
\end{gather*}

\subproblem{c}
\begin{gather*}
    K_n = \int_0^{\frac{\pi}{2}} \cos^n x \sin (nx)\, dx\\
    =\begin{bmatrix}
        u=\cos^n x & dv=\sin (nx)\, dx\\
        du=-n\cos^{n-1}x\sin x  & v=-\frac{1}{n}\cos (nx)
    \end{bmatrix} -\frac{1}{n}\cos^n x\cos nx\Big|^\frac{\pi}{2}_0 \\
    + \int_0^\frac{\pi}{2} \cos ^{n-1}x\sin x \cos (nx)\, dx = \frac{1}{n} - \int_0^\frac{\pi}{2} \cos ^{n-1}x\sin x \cos (nx)\, dx\\
    \sin((n-1)x) = \sin nx \cos x - \cos nx \sin x\\
    K_{n-1} = \int_0^{\frac{\pi}{2}} \cos^{n-1}x \sin((n-1)x)\, dx \\= \int_0^{\frac{\pi}{2}} \cos^{n-1}x (\sin nx \cos x - \cos nx \sin x)\, dx\\
    = \int_0^{\frac{\pi}{2}} \cos^n x \sin nx \, dx - \int_0^{\frac{\pi}{2}} \cos^{n-1} x \sin x \cos nx \, dx  = 2K_n - \frac{1}{n}\\
    \boxed{K_n = \frac{1}{2}K_{n-1} + \frac{1}{2n}}\\
    K_n = \frac{1}{2}K_{n-1} + \frac{1}{2n} = \frac{1}{4}K_{n-2} + \frac{1}{2n} + \frac{1}{4(n-1)}\\
    = \dots = \paren{\frac{1}{2}}^nK_0 + \sum_{k=1}^n \paren{\frac{1}{2}}^{n-k+1}\frac{1}{k}\\
    \boxed{K_n = \frac{1}{2^{n+1}}\sum_{k=1}^n \frac{2^k}{k}}
\end{gather*}

\subproblem{d}
\begin{gather*}
    I_1=\frac{\pi}{4}, I_2=\frac{\pi}{8}\\
    K_1=\frac{1}{2},K_2=\frac{1}{2}
\end{gather*}


























\end{document}

